\chapter{Panduan Asisten Praktikum}

\section{Peran Asisten}

Asisten memfasilitasi proses belajar tanpa mengambil alih pekerjaan mahasiswa.
Seluruh DDL, query, skrip ETL, analisis, dan laporan tetap dikerjakan sendiri
oleh mahasiswa. Asisten bertanggung jawab untuk:

\begin{enumerate}[leftmargin=1.4em]
  \item memverifikasi kesiapan lingkungan sebelum sesi dimulai;
  \item menghubungkan konsep kuliah dengan objek basis data yang dibangun di
        laboratorium;
  \item menjalankan checkpoint pada waktu yang ditentukan menggunakan skrip
        pemeriksa;
  \item membantu diagnosis berdasarkan pesan galat dan keluaran
        \perintah{EXPLAIN}, bukan berdasarkan tebakan;
  \item menjaga agar data mentah pada basis data sumber tidak diubah; dan
  \item menilai menggunakan rubrik yang sama untuk seluruh kelas.
\end{enumerate}

\section{Alur Sesi 120 Menit}

Karena waktu tatap muka terbatas, pemasangan perangkat lunak dan pembacaan
teori dasar harus sudah diselesaikan mahasiswa sebelum sesi. Asisten memakai
waktu kelas untuk menghubungkan teori dengan objek yang dibangun, memandu
prosedur inti, dan memastikan setiap mahasiswa mencapai checkpoint minimal.

\begin{center}
\small
\begin{tabular}{@{}p{0.62\textwidth}r@{}}
\toprule
\textbf{Aktivitas} & \textbf{Durasi} \\
\midrule
Briefing dan demonstrasi satu contoh oleh asisten & 15 menit \\
Kerja terbimbing pada lembar kerja, asisten berkeliling & 80 menit \\
Checkpoint dengan skrip pemeriksa otomatis & 15 menit \\
Briefing tugas luar laboratorium & 10 menit \\
\midrule
\textbf{Total} & \textbf{120 menit} \\
\bottomrule
\end{tabular}
\end{center}

\section{Target Waktu dan Checkpoint}

\begin{itemize}
  \item Menit 0--15: mahasiswa memahami kasus, luaran yang dituju, dan satu
        contoh yang dikerjakan asisten di depan kelas.
  \item Menit 15--95: mahasiswa menyelesaikan Bagian A sampai D pada prosedur
        terbimbing dan menghasilkan artefak yang diminta modul.
  \item Menit 95--110: asisten menjalankan skrip pemeriksa dan menangani butir
        yang gagal. Skrip mencetak lulus atau gagal per butir sehingga asisten
        hanya perlu menangani yang gagal --- inilah yang membuat 15 menit cukup
        untuk satu kelas.
  \item Menit 110--120: asisten menjelaskan tugas luar laboratorium, format
        berkas, dan tenggat pengumpulan.
\end{itemize}

Jika demonstrasi mengalami kendala teknis, asisten tidak boleh mengorbankan
waktu kerja terbimbing. Gunakan \berkas{snapshot.sql} modul sebelumnya yang
telah disiapkan, lalu lanjutkan diagnosis setelah sesi.

\section{Batas Bantuan}

Asisten boleh memberi petunjuk, menunjukkan dokumentasi, membantu membaca pesan
galat PostgreSQL, dan membantu menafsirkan keluaran \perintah{EXPLAIN}. Asisten
tidak memberikan salinan DDL atau query jawaban, tidak memperbaiki seluruh
skrip mahasiswa, dan tidak menentukan kesimpulan analisis.

\begin{catatanasisten}
Kaidah yang berlaku sepanjang praktikum: \textbf{setiap keputusan teknis harus
dapat dijelaskan}. Ketika mahasiswa bertanya ``tipe data apa yang benar?'',
pertanyaan balik yang tepat adalah ``berapa nilai terbesar yang mungkin, dan
berapa byte yang Anda bersedia bayar untuk itu?''. Jawaban benar tanpa alasan
memperoleh separuh nilai.
\end{catatanasisten}

\section{Checklist Sebelum Sesi}

\begin{ceklis}
  \item \berkas{docker compose up} telah diuji pada Windows, macOS, dan Linux.
  \item Basis data sumber dan gudang dapat diakses dari jaringan laboratorium.
  \item Dataset berukuran sesuai kebutuhan modul telah tersedia --- ukuran
        kecil untuk Modul 1--5, ukuran sedang untuk Modul 6--8.
  \item Nilai \perintah{--seed} pembangkit data sama untuk seluruh kelas
        sehingga angka hasil profiling dapat dibandingkan antarkelompok.
  \item \berkas{snapshot.sql} keadaan akhir modul sebelumnya tersedia bagi
        mahasiswa yang tertinggal.
  \item Skrip pemeriksa \berkas{check.sql} atau \berkas{check.py} telah
        dijalankan pada basis data referensi dan seluruh butirnya lulus.
  \item Ekstensi PostgreSQL yang diperlukan modul telah diaktifkan --- lihat
        Tabel~\ref{tab:ekstensi}.
  \item Rubrik dan format pengumpulan tersedia bagi mahasiswa.
\end{ceklis}
